\documentclass[11pt]{article}

% ---------------------------------------------------------------------------
% Generic scaffold. Swap this preamble for a venue's official class/style
% file once a target is chosen (e.g. arXiv, Quantum Machine Intelligence,
% an IEEE QCE workshop template) -- everything below \begin{document} should
% survive that swap mostly unchanged.
% ---------------------------------------------------------------------------
\usepackage[margin=1in]{geometry}
\usepackage{amsmath,amssymb}
\usepackage{graphicx}
\usepackage{booktabs}
\usepackage{natbib}
\usepackage{hyperref}
\usepackage{xcolor}
\usepackage{subcaption}

% Figures are read directly from the results directories that produced them
% (no duplicated binaries in paper/) -- add new run directories here as needed.
\graphicspath{%
  {figures/}%
  {../qcpinn_results/20260712_035925/}%
  {../qcpinn_results/ablation_20260724_020639/}%
  {../qcpinn_results/ablation2_20260729_225346/}%
  {../pinn_results/20260625_161036/}%
}

\newcommand{\todo}[1]{\textcolor{red}{[TODO: #1]}}

\title{Do Variational Quantum Circuits Help Physics-Informed Neural Networks?\\
A Systematic Ablation Study on a Diffusion PDE Benchmark}
% TODO: finalize title with advisor -- see paper/README.md for alternatives.

\author{\todo{author list} \\ \todo{affiliation}}
\date{}

\begin{document}
\maketitle

\begin{abstract}
% Draft abstract -- rewrite once Discussion/Conclusion are final; abstracts
% should be written last but a placeholder here keeps main.tex compiling.
Physics-informed neural networks (PINNs) that replace part of the network
with a variational quantum circuit (QCPINN) have recently been reported to
match or exceed classical PINNs while using far fewer trainable parameters.
We evaluate this claim on a seawater temperature diffusion PDE with three
boundary-condition types (Dirichlet, Neumann, Robin) and both forward and
inverse problem variants, using the quantum circuit ansatz proposed by
\citet{farea2025qcpinn}. Across all six scenarios, the quantum branch is
both less accurate (\todo{X--Y}$\times$ higher mean absolute error) and
substantially slower in wall-clock time (\todo{X--Y}$\times$) than the
classical baseline. Through a two-stage ablation -- first over training
hyperparameters (initialization scale, learning-rate schedule), then over
circuit structure (qubit count, circuit depth) with direct measurement of
the circuit-only gradient norm -- we rule out hyperparameter choice as the
cause and find a structure-dependent effect that fits neither a clean
expressivity-limited nor a clean barren-plateau account. We discuss why
wall-clock accounting, largely absent from prior QCPINN evaluations, changes
the practical picture, and what our gradient measurements add to existing
theoretical barren-plateau claims.

\end{abstract}

\section{Introduction}\label{sec:intro}
% TODO(zh): 这节需要你自己的语言组织，下面给出段落级骨架和每段该完成的任务。

% Paragraph 1 -- PINN background and motivation for the specific PDE.
Physics-informed neural networks \citep{raissi2019pinn} solve forward and
inverse partial differential equation (PDE) problems by embedding the
governing equation, boundary conditions, and (for inverse problems)
observational data directly into the training loss of a neural network
surrogate. \todo{One or two sentences motivating the seawater temperature
diffusion equation specifically -- cite the paper you are replicating/
extending (Han et al., 2026, Journal of Oceanology and Limnology) and state
why it is a useful benchmark: closed-form analytical solutions exist for
all three boundary-condition types, which makes ground-truth error
measurable exactly rather than via a numerical reference solution.}

% Paragraph 2 -- QML hype + hybrid quantum-classical PINNs as a response.
\todo{Motivate interest in quantum machine learning for scientific
computing in one or two sentences, then introduce the hybrid
quantum-classical PINN (QCPINN/QPINN) line of work as the object of study.
Keep this short -- the literature survey belongs in Related Work, not here.}

% Paragraph 3 -- the gap / research questions.
Despite growing enthusiasm for hybrid quantum-classical PINNs, we find that
efficiency in this literature is reported almost exclusively via parameter
count or number of training epochs, rarely via wall-clock time -- a
distinction that matters because variational quantum circuits are, in the
near term, executed via classical simulation whose per-step cost scales
with circuit width and depth (see Section~\ref{sec:related}). This raises
two questions that motivate the present study:
\begin{itemize}
  \item[\textbf{RQ1}] On a benchmark PDE with exact analytical solutions,
    does a QCPINN built on a previously proposed circuit ansatz
    \citep{farea2025qcpinn} match classical PINN accuracy once wall-clock
    training time is accounted for, not just parameter count?
  \item[\textbf{RQ2}] If it does not, is the gap attributable to
    training hyperparameters (initialization, learning-rate schedule), or
    to the quantum circuit's structure (expressivity or trainability)?
\end{itemize}

% Paragraph 4 -- contributions, stated as a list.
We make the following contributions:
\begin{itemize}
  \item A controlled classical-vs-quantum comparison across six matched
    scenarios (three boundary-condition types $\times$ forward/inverse),
    with explicit wall-clock accounting alongside accuracy metrics
    (Section~\ref{sec:experiments}).
  \item A two-phase ablation that first rules out training hyperparameters
    as the cause of the quantum branch's underperformance, then isolates
    circuit structure (qubit count, depth) while directly measuring the
    circuit-only gradient norm -- a diagnostic largely absent from prior
    QCPINN evaluations, which infer trainability issues indirectly from
    accuracy curves or rely on purely theoretical gradient-decay bounds
    (Section~\ref{sec:experiments}).
  \item A discussion of why our findings diverge from prior QCPINN results
    obtained with a closely related circuit design, and what that implies
    for how QCPINN efficiency should be reported going forward
    (Section~\ref{sec:discussion}).
\end{itemize}

\todo{Close with a one-sentence roadmap of the paper's structure once the
section numbers/titles are final.}


\section{Related Work}\label{sec:related}
% Revised 2026-08-28 after full-text verification of farea2025qcpinn and
% hqpinn2025flow (see references.bib changelog comment). Numbers and quotes
% below are checked against the source PDFs, not abstract summaries.
% Remaining citations (qcpinn2025reservoir, hqcpinn2026hydro, qpinnmac2025)
% are still unverified stubs -- see references.bib TODOs.

Physics-informed neural networks (PINNs) \citep{raissi2019pinn} have become
a standard tool for solving forward and inverse PDE problems by embedding
the governing equations directly into the training loss. Motivated by
broader interest in quantum machine learning for scientific computing, a
growing line of work replaces part of the classical network with a
variational quantum circuit (VQC), typically in a hybrid architecture where
classical layers handle pre- and post-processing while the VQC provides the
trainable core -- an approach variously termed QPINN or QCPINN.

\citet{farea2025qcpinn} introduce QCPINN and evaluate it, alongside a
continuous-variable (qumode-based) variant that trains very poorly across
the board, on five forward-only PDE benchmarks with Dirichlet-type boundary
conditions: Helmholtz, a 2D lid-driven cavity flow, a 1D wave equation,
Klein-Gordon, and convection-diffusion. Their headline claim -- 10--30\% of
the trainable parameters of a classical PINN with comparable or better
accuracy -- holds for three of the five: a 4--64\% reduction in relative
$L_2$ error on Helmholtz, Klein-Gordon, and convection-diffusion. On the
remaining two, the quantum branch is \emph{less} accurate than the
classical baseline despite using far fewer parameters (cavity: 11--23\%
higher relative error across the three output fields; wave: 12\% higher),
which the authors state plainly rather than omit \citep[Sec.~8,
``Conclusion'']{farea2025qcpinn}. They also report a timing column absent
from the abstract: per-iteration wall-clock time is 15--40$\times$ that of
the classical baseline in every one of the five benchmarks, which they
attribute to the cost of simulating the circuit classically rather than to
the algorithm itself. Separately, their qubit-based circuit sweep tests
four named topologies (Alternate, Cascade, Cross-mesh, Layered) crossed
with two data-encoding schemes (amplitude, angle); the accuracy and timing
numbers above are for each benchmark's best-performing topology, most
often Cascade or Cross-mesh -- both of which use continuously-tunable
controlled-rotation entanglers (CRX/CRZ) rather than fixed CNOT gates,
which the authors argue is why those two topologies converge more reliably
than Alternate or Layered.

\citet{hqpinn2025flow} report a closely related pattern on a different
problem family: the compressible Euler equations, tested across a smooth
harmonic traveling wave, a non-harmonic solution with a moving contact
discontinuity, and a 2D transonic aerofoil flow with a shock. A quantum
branch whose functional form is an explicit truncated Fourier series
achieves the single best loss of any model on the harmonic case (an
order-of-magnitude improvement over the best classical network, using far
fewer parameters), but ``significantly underperformed'' on the non-harmonic
case and underperformed further still on the shock-dominated transonic
case, where the quantum branch is additionally under-parameterized. The
paper's own stated thesis is that harmonicity of the target solution, not
problem difficulty in general, is what determines whether the quantum
branch helps: ``quantum PINNs excel in harmonic regimes... due to the
Fourier structure of PQCs... [but] struggle to generalize to
discontinuities or shocks.'' On timing, wall-clock cost is reported only
for the transonic case, where per-epoch cost for the quantum configurations
ranges 17--82$\times$ that of the classical configurations tested (the
paper's own summary rounds this to ``up to 100$\times$''); a separate,
larger ``two orders of magnitude'' figure appearing earlier in the paper is
attributed to other authors' measurements, not to this paper's own
experiments. Notably, this paper does not use the term ``barren plateau''
anywhere; its account of trainability loss in larger quantum models is
framed instead as ``excessive expressibility,'' a distinct (if related)
diagnosis.

Read together, these two papers already show that the value of a quantum
branch in a hybrid PINN is problem-dependent rather than uniform -- within
a single paper's own benchmark suite, in both cases -- and that reporting
practices vary: parameter-count efficiency is reported everywhere, but
wall-clock or per-iteration time is reported only for a subset of
benchmarks in each paper (all five in \citealp{farea2025qcpinn}, only the
hardest of three in \citealp{hqpinn2025flow}), which limits how directly
either paper's timing claims generalize.
\todo{Add hydrological \citep{hqcpinn2026hydro}, reservoir seepage
\citep{qcpinn2025reservoir}, and the theoretical barren-plateau bound
\citep{qpinnmac2025} once those are verified against full text -- see
references.bib TODOs. Until then, do not restate the epoch-count /
polynomial-gradient-bound claims for these three beyond what the earlier
abstract-level search already summarized in chat, since that summary has
not been checked against the full papers the way the two above have.}

Our work differs from this pair of closely related studies in three ways.
First, both restrict the quantum branch to forward, Dirichlet-type problems
(\citealp{farea2025qcpinn} exclusively; \citealp{hqpinn2025flow} uses
Dirichlet conditions throughout as well); we test both forward and inverse
formulations across three boundary-condition types (Dirichlet, Neumann,
Robin) on a single diffusion PDE with closed-form solutions in all six
cases, which lets us measure error against ground truth exactly rather than
against a numerical or experimental reference. Second, where
\citet{farea2025qcpinn} report per-iteration time as one column among many
and \citet{hqpinn2025flow} report it for only one of three benchmarks, we
report wall-clock training time as a primary outcome for every scenario
alongside accuracy, and find the quantum branch uniformly worse on both
axes in all six -- a stronger and more consistent negative result than
either prior paper's own (partially mixed) findings. Third, neither prior
paper isolates circuit structure as an experimental variable with direct
gradient measurement: \citet{farea2025qcpinn} discusses trainability
differences between entangler choices (CRX/CRZ vs.\ fixed CNOT)
qualitatively, and formally invokes barren plateaus only for its
continuous-variable circuit, which is not the qubit-based design we build
on; \citet{hqpinn2025flow} attributes reduced trainability in larger
circuits to ``excessive expressibility'' based on final-loss curves, not a
direct gradient-norm measurement. We instrument the circuit-only gradient
norm directly across a controlled qubit-count/depth sweep, providing an
empirical data point neither paper's diagnostic approach supplies.


\section{Method}\label{sec:method}
\subsection{Problem setup}
We consider the one-dimensional temperature diffusion equation
\begin{equation}
  \frac{\partial T}{\partial t} = k_z \frac{\partial^2 T}{\partial z^2},
  \qquad z \in [0, 2],\ t \in [0, 5],
\end{equation}
with true diffusivity $k_z = 0.05$, following \citet{han2026seawater}.
% NOTE: han2026seawater in references.bib is still an empty stub -- I
% cannot fabricate its author list, volume, or page numbers. This is the
% one remaining reference that has to come from you/your advisor rather
% than from a web search, since it's your own group's source paper.
We evaluate three boundary
conditions, each with a closed-form analytical solution used as ground
truth:
\begin{itemize}
  \item \textbf{Dirichlet}: $T = e^{-k_z(\pi/2)^2 t}\sin(\pi z/2) +
    e^{-k_z \pi^2 t}\sin(\pi z)$.
  \item \textbf{Neumann}: the corresponding cosine-series solution.
  \item \textbf{Robin}: $\partial T/\partial z + T = 0$ at $z=0,2$; the
    general-term solution follows \citet{han2026seawater}, Eq.~18.
\end{itemize}
For each boundary condition we solve both a \emph{forward} problem
($k_z$ known, solve for $T$) and an \emph{inverse} problem ($k_z$ unknown,
estimated jointly with $T$ from sparse observations), giving six scenarios
in total.

\subsection{Classical PINN}
The classical baseline is a fully connected network mapping $(z, t) \mapsto
T$: five hidden layers of 32 units with $\tanh$ activations. The loss
combines the PDE residual, boundary-condition residual, and (for inverse
problems) a data-fit term against sparse noisy observations; $k_z$ is a
trainable scalar in the inverse case (parameterized as $\exp(\log k_z)$ for
positivity). Training samples 2540 PDE-residual (collocation) points, 80
boundary points (40 per side), and 160 initial-condition points, each
redrawn uniformly at random per run; inverse scenarios additionally use 12
noiseless observation points on a fixed $4\times 3$ $(z,t)$ grid. The loss
term weights are $w_{bc}=w_{ic}=w_{data}=1$ throughout, and $w_{pde}=1$ for
Dirichlet and Neumann scenarios; Robin scenarios use $w_{pde}=10$, since
the gradient-dependent Robin boundary condition is harder to satisfy at
$w_{pde}=1$ and needs the PDE residual up-weighted to converge (an
asymmetry inherited from the original classical replication rather than
introduced for this study). Training is two-phase: Adam ($\mathrm{lr}=
10^{-3}$) for a fixed number of steps (5000 for Dirichlet/Neumann, 20000
for Robin, matching the higher $w_{pde}$), followed by L-BFGS (strong-Wolfe
line search, history size 100, $\mathrm{tolerance\_grad}=10^{-9}$,
$\mathrm{tolerance\_change}=10^{-11}$, up to 10000 or 20000 iterations
matching the Adam budget) until L-BFGS's own convergence tolerances are
met or the iteration budget is exhausted, whichever comes first -- there
is no early stopping based on held-out validation error.

\subsection{QCPINN}
The quantum branch is inspired by, but not identical to, the qubit-based
(discrete-variable) circuit family in \citet{farea2025qcpinn}. Inputs
$(z,t)$ are linearly rescaled to $[-1,1]$ and encoded via
\texttt{AngleEmbedding} (RX rotations) onto $n$ qubits, matching their
angle-embedding scheme. Each of $L$ variational layers applies RZ-RX
rotations on every qubit, a \emph{ring} of CNOT entangling gates (qubit $i$
to qubit $(i+1) \bmod n$, wrapping from the last qubit back to the first),
then a mirrored RX-RZ rotation layer; the circuit output is the vector of
Pauli-$Z$ expectation values $\langle Z_i\rangle \in [-1,1]$, matching
their readout. \citet{farea2025qcpinn} test four named topologies, none of
which is exactly this combination: their \emph{Cascade} topology uses ring
connectivity but continuously-tunable CRX gates rather than fixed CNOT,
while their \emph{Alternate} and \emph{Layered} topologies use fixed CNOT
but nearest-neighbor (not ring) connectivity. We therefore refer to our
circuit as \emph{ring-CNOT}: a fifth ring-plus-fixed-entangler combination
that is closest in spirit to their Cascade topology but was not itself
among the four they swept, rather than an instance of any one of their
named topologies. This output is post-processed by a small
classical head to produce $T$ (and, for inverse problems, $k_z$). Circuit
parameters are simulated exactly (no shot noise, no hardware noise) via
PennyLane's \texttt{default.qubit} device with the \texttt{backprop}
differentiation method -- \citet{farea2025qcpinn} state they use the same
shots-None/backprop configuration for their qubit-based circuits but do not
print the specific simulator device string, so this match is inferred, not
confirmed verbatim. The baseline configuration uses $n=4$ qubits and $L=2$
layers; note this differs from the $n=5$, $L=1$ configuration
\citet{farea2025qcpinn} found best after their own sweep -- Section
\ref{sec:experiments} sweeps qubit count and depth independently around our
own baseline, not theirs. The quantum branch is trained differently from
the classical one, not only architecturally but in optimizer choice: it
never switches to L-BFGS -- each L-BFGS step would require multiple
closure evaluations through the quantum circuit, which is too expensive
per step -- and instead runs Adam throughout, for a step budget equal to
the classical branch's combined Adam+L-BFGS budget (i.e.\ 15000 steps for
Dirichlet/Neumann, 40000 for Robin), with cosine learning-rate decay from
$10^{-3}$ to a floor $\eta_{\min}=10^{-5}$. This means the two branches
differ in optimizer (first-order throughout vs.\ first-order then
second-order) as well as in network type, a confound Section
\ref{sec:experiments}'s phase-1 ablation partially addresses by varying
$\eta_{\min}$ but does not eliminate -- both branches were given a matched
total step/iteration budget, not a matched optimizer.

\subsection{Training and evaluation}
All experiments use PyTorch with fixed seeds. Seeding differs by
experiment: the main classical-vs-quantum comparison
(Section~\ref{sec:experiments}, Table~\ref{tab:main-results}) uses a
single seed (0) per scenario, so the reported numbers there are point
estimates with no variance across runs; the phase-1 and phase-2 ablations
each use 3 seeds (0, 1, 2) per (scenario, configuration) pair, and it is
only for those that we report a mean and standard deviation. We return to
what this asymmetry means for how the main comparison should be read in
Section~\ref{sec:discussion}'s limitations. We report:
\begin{itemize}
  \item \textbf{MAE} and \textbf{relative $L_2$ error} against the
    analytical solution on a held-out grid;
  \item \textbf{PDE residual} at collocation points, as an
    optimization-quality check independent of the analytical solution;
  \item \textbf{$k_z$ estimation error} for inverse scenarios;
  \item \textbf{Wall-clock training time}, measured on an Apple M2 (8
    cores, 16GB RAM), CPU only for both branches -- neither PyTorch nor
    PennyLane was given GPU access, so the reported ratios are a CPU-only
    comparison. This likely does not favor the classical branch
    disproportionately: a GPU-accelerated classical network would train
    faster still, which would widen rather than narrow the gap, whereas
    the quantum branch's cost is dominated by classical simulation of the
    circuit itself and would not benefit comparably from the same
    hardware upgrade;
  \item \textbf{Circuit-only gradient norm} (phase-2 ablation only): the
    norm of the loss gradient with respect to the VQC parameters alone,
    logged at the start and end of training, isolating circuit
    trainability from the classical pre/post-processing layers.
\end{itemize}


\section{Experiments}\label{sec:experiments}
\subsection{Classical vs.\ quantum: main comparison}
Table~\ref{tab:main-results} compares the classical and quantum branches
across all six scenarios. The quantum branch is less accurate in every
scenario (1.5--14.4$\times$ higher MAE) and substantially slower in every
scenario (53--89$\times$ wall-clock time), with no scenario favoring the
quantum branch on either axis.
% Table auto-generated from qcpinn_results/20260712_035925/ -- rerun
% tables/make_tables.py after any change to that CSV.
% Auto-generated by tables/make_tables.py -- do not edit by hand.
\begin{table}[t]
\centering
\caption{Classical vs.\ quantum PINN across six boundary-condition / problem-type
scenarios. MAE and wall-clock time (s) are means over seeds; ratio columns are
quantum relative to classical (higher = quantum worse).}
\label{tab:main-results}
\begin{tabular}{llrrrrrr}
\toprule
BC & Problem & MAE (cls.) & MAE (qtm.) & MAE ratio & Time (cls.) & Time (qtm.) & Time ratio \\
\midrule
Dirichlet & Forward & 0.000231 & 0.00034 & 1.5$\times$ & 31.9 & 2.44e+03 & 76.5$\times$ \\
Dirichlet & Inverse & 0.000287 & 0.0005 & 1.7$\times$ & 39.1 & 2.69e+03 & 69.0$\times$ \\
Neumann & Forward & 0.000166 & 0.00239 & 14.4$\times$ & 38.4 & 2.56e+03 & 66.6$\times$ \\
Neumann & Inverse & 0.000249 & 0.00118 & 4.7$\times$ & 35.6 & 2.82e+03 & 79.3$\times$ \\
Robin & Forward & 8.96e-05 & 0.000384 & 4.3$\times$ & 128 & 6.79e+03 & 52.9$\times$ \\
Robin & Inverse & 8.66e-05 & 0.00029 & 3.3$\times$ & 133 & 1.19e+04 & 89.3$\times$ \\
\bottomrule
\end{tabular}
\end{table}


\todo{Add 1-2 representative solution/error plots here, e.g.
\texttt{DE\_Dir\_FWD\_classical.png} vs \texttt{DE\_Dir\_FWD\_quantum.png}
from qcpinn\_results/20260712\_035925/ (already on the graphicspath in
main.tex), to show what the accuracy gap looks like spatially, not just in
the aggregate MAE number.}

\subsection{Ablation 1: training hyperparameters}
The main comparison leaves open whether the quantum branch's
underperformance reflects a suboptimal training recipe rather than a
structural limitation. We test two candidates on two representative
scenarios (Dirichlet-forward, Robin-inverse): a smaller weight
initialization scale (\texttt{small\_init}) intended to keep VQC rotation
angles away from regions with vanishing gradient, and a higher
learning-rate floor (\texttt{slow\_decay}) intended to prevent premature
convergence to a poor optimum. Table~\ref{tab:ablation1} shows neither
closes the gap: \texttt{slow\_decay} is the worst of the three quantum
configurations in both scenarios, and \texttt{small\_init} yields only a
3--16\% MAE improvement over the untouched baseline -- far short of the
3--14$\times$ gap to the classical branch.
% Auto-generated by tables/make_tables.py -- do not edit by hand.
\begin{table}[t]
\centering
\caption{Phase-1 ablation: quantum-branch initialization scale and LR-schedule
floor, vs.\ the classical reference and the untouched quantum baseline.
Neither intervention closes the accuracy or runtime gap.}
\label{tab:ablation1}
\begin{tabular}{llrr}
\toprule
Scenario & Config & MAE (mean $\pm$ std) & Time (s) \\
\midrule
Dirichlet Forward & classical & 0.000171 $\pm$ 3.7e-05 & 65.68 \\
Dirichlet Forward & baseline & 0.000522 $\pm$ 0.00013 & 5257 \\
Dirichlet Forward & small\_init & 0.000437 $\pm$ 6e-05 & 5097 \\
Dirichlet Forward & slow\_decay & 0.00102 $\pm$ 0.0008 & 5038 \\
Robin Inverse & classical & 7.91e-05 $\pm$ 2e-05 & 400.3 \\
Robin Inverse & baseline & 0.000626 $\pm$ 0.00033 & 1.492e+04 \\
Robin Inverse & small\_init & 0.000609 $\pm$ 0.00023 & 1.513e+04 \\
Robin Inverse & slow\_decay & 0.0011 $\pm$ 0.00032 & 1.476e+04 \\
\bottomrule
\end{tabular}
\end{table}


\subsection{Ablation 2: circuit structure}
Having ruled out these two hyperparameters, we next ask whether the circuit
itself is the bottleneck -- either under-expressive (too few qubits/layers)
or subject to a barren-plateau-like effect (gradients vanishing as the
circuit grows). We sweep qubit count and depth independently around the
baseline ($n{=}4, L{=}2$): \texttt{narrow} ($n{=}2$), \texttt{wide}
($n{=}6$), \texttt{deep} ($L{=}4$), holding initialization and
learning-rate schedule at the phase-1 baseline throughout. Alongside the
usual accuracy metrics we log the circuit-only gradient norm at the end of
training. Table~\ref{tab:ablation2} summarizes the result.
% Auto-generated by tables/make_tables.py -- do not edit by hand.
\begin{table}[t]
\centering
\caption{Phase-2 ablation: circuit width/depth sweep (narrow/baseline/wide/deep)
with final circuit-only gradient norm. Accuracy and gradient norm do not move
monotonically together, so neither a pure expressivity nor a pure
barren-plateau account fits cleanly -- see Discussion.}
\label{tab:ablation2}
\begin{tabular}{llrrr}
\toprule
Scenario & Config & MAE (mean $\pm$ std) & Final $\|\nabla_{\text{circuit}}\|$ & Time (s) \\
\midrule
Dirichlet Forward & classical & 0.000184 $\pm$ 3.5e-05 & 0 & 31.96 \\
Dirichlet Forward & baseline & 0.00052 $\pm$ 0.00013 & 1.56e-05 & 2655 \\
Dirichlet Forward & narrow & 0.00152 $\pm$ 0.00052 & 1.93e-05 & 968.1 \\
Dirichlet Forward & wide & 0.000446 $\pm$ 8.9e-05 & 9.45e-06 & 7398 \\
Dirichlet Forward & deep & 0.000663 $\pm$ 0.00023 & 2.21e-05 & 4431 \\
Robin Inverse & classical & 6.44e-05 $\pm$ 1.6e-05 & 0 & 151.3 \\
Robin Inverse & baseline & 0.000627 $\pm$ 0.00033 & 0.00653 & 8331 \\
Robin Inverse & narrow & 0.00177 $\pm$ 0.00092 & 0.0326 & 3173 \\
Robin Inverse & wide & 0.000413 $\pm$ 0.00013 & 0.00575 & 2.201e+04 \\
Robin Inverse & deep & 0.000508 $\pm$ 0.0001 & 0.00595 & 1.393e+04 \\
\bottomrule
\end{tabular}
\end{table}


\todo{Add the ablation2\_accuracy.png / ablation2\_qgrad.png figures from
qcpinn\_results/ablation2\_20260729\_225346/ here as a side-by-side, since
the table alone makes the non-monotonic pattern harder to see than the
plots do.}

The pattern is mixed rather than clean. \texttt{wide} achieves the best
accuracy among the quantum configurations in both scenarios, which would
suggest expressivity is a limiting factor worth scaling past -- but it also
has the \emph{smallest} final gradient norm of the four configurations,
which is not the direction a pure expressivity story predicts and is only
partially consistent with a barren-plateau account (which predicts smaller
gradients accompanying \emph{worse}, not better, accuracy). \texttt{narrow}
is the worst-performing configuration on accuracy, consistent with an
under-expressivity explanation at the low end, while \texttt{deep} is worse
than baseline on accuracy despite a comparable gradient norm. \todo{State
explicitly in the text -- and defend in Discussion -- that this rules out a
single clean explanation and motivates the more cautious framing adopted
there (data encoding / ansatz choice as an unresolved third factor).}


\section{Discussion}\label{sec:discussion}
% Revised 2026-08-28 after full-text verification of farea2025qcpinn and
% hqpinn2025flow. Key corrections vs. the earlier abstract-only draft:
% (1) farea2025qcpinn DOES report per-iteration timing (15-40x), just not
%     in the abstract, and its own results are already mixed (worse on 2/5
%     PDEs) -- don't overstate it as a uniformly positive prior result;
% (2) hqpinn2025flow's own measured max ratio is ~82x, not a clean 100x --
%     that rounded figure is their own summary of their own Table A3, while
%     a separate "two orders of magnitude" claim earlier in their paper is
%     attributed to OTHER authors, not this paper's own experiment;
% (3) "barren plateau" is NOT a term used in hqpinn2025flow at all (they
%     say "trainability" / "excessive expressibility" instead); the one
%     paper that does invoke barren plateaus explicitly is
%     farea2025qcpinn's continuous-variable circuit discussion -- which is
%     a DIFFERENT circuit paradigm from the qubit-based one we build on,
%     so don't cite that barren-plateau discussion as if it were about the
%     same architecture family.

\subsection{Prior positive results are narrower than they first appear}
None of \citet{farea2025qcpinn}, \citet{hqpinn2025flow}, or
\citet{qcpinn2025reservoir} is a uniform "quantum wins" result once read
past the abstract: all three contain within-paper negative or mixed cases
-- cavity flow and the wave equation for the first, non-harmonic and
transonic regimes for the second, no circuit topology winning across all
four benchmarks (plus an explicit disclaimer against causal claims) for
the third. We read our six-scenario result not as contradicting these
papers but as a more consistent, more uniform version of a pattern they
already show in part: where they find a mix depending on PDE, regime, or
topology, we find the quantum branch losing on every single one of six
matched scenarios on a problem family -- a diffusion equation under three
boundary-condition types, in both forward and inverse form -- that none of
the three tests at all: all three restrict to forward problems with
Dirichlet (and, in one case, additionally Neumann) conditions, and none
includes an inverse/parameter-estimation variant. It is also worth being
precise about how much independent weight \citet{qcpinn2025reservoir}
carries: it is not an arms-length third data point, since it explicitly
reuses \citeauthor{farea2025qcpinn}'s open-source implementation for a new
PDE domain. Its convergence with \citet{hqpinn2025flow} on the qualitative
claim that a quantum branch's advantage is regime- or topology-dependent
rather than uniform therefore reflects two independent author teams
agreeing on that shape, not three independent confirmations of any
specific number.

\subsection{Wall-clock accounting: sharper than the literature, but comparable}
\citet{farea2025qcpinn} report per-iteration time (15--40$\times$ slower
than classical, across all five of their benchmarks) and
\citet{hqpinn2025flow} report per-epoch time for one of their three
benchmarks (17--82$\times$, which they themselves round to "up to
100$\times$"). \citet{qcpinn2025reservoir} and \citet{hqcpinn2026hydro} sit
at the far end of this spectrum: the former runs on a fully specified CPU
(two Intel Xeon E5-2680 v4, 28 cores, 128GB RAM) and explicitly
acknowledges in prose that "quantum simulation has inherent overhead on
classical hardware," yet reports no timing measurement whatsoever; the
latter reports an epoch-count "speedup" (fewer epochs to reach a
validation-loss threshold) without ever mentioning wall-clock cost at all,
despite using parameter-shift-rule gradients -- known to require multiple
extra circuit evaluations per step relative to backprop -- which makes an
epoch-count comparison an especially weak proxy for real training cost in
that specific case. Two of four closely related papers reporting zero
timing data, and a third reporting it for only one of three benchmarks, is
itself worth stating as a field-wide pattern rather than a coincidence
across independent author teams. Our own 53--89$\times$ wall-clock ratio is
in the same direction as the two papers that do measure it, and a similar
order of magnitude -- not a qualitatively new finding -- but we measure it
as a primary outcome across all six of our scenarios rather than as a
secondary column reported for some subset of benchmarks, or omitted
entirely. Our ratio runs somewhat higher than Farea et al.'s 15--40$\times$;
one plausible source is circuit size, since their best-performing
configuration uses only 5 qubits and 1 layer, smaller than our 4-qubit,
2-layer baseline (and our own ablation sweeps up to 6 qubits and 4 layers,
which should be commensurately slower still). However, none of the three
papers states enough hardware detail to rule out a hardware difference as
a contributing factor instead or in addition -- Farea et al.\ report a
6148-CPU, 192GB cluster node; the flow paper reports only "distributed
CPU, no GPU"; the reservoir paper specifies its CPU but has no timing to
compare against at all -- so this is a limit on how precisely wall-clock
ratios can be compared \emph{across} papers, not only a feature of our own
results.

\subsection{A smooth problem that still loses -- and why "barren plateau" needs care}
\citet{hqpinn2025flow}'s central thesis is that harmonicity of the target
solution predicts whether the quantum branch helps, not problem difficulty
in general. Our diffusion PDE has closed-form sine/cosine solutions and so
sits on the "harmonic" side of their distinction, yet our quantum branch
loses in all six scenarios, including the smoothest ones. This is a
genuine point of tension worth naming directly: either the
harmonic/non-harmonic axis from their compressible-flow, Fourier-series
argument does not transfer to a diffusion PDE encoded via AngleEmbedding,
or some other factor in our setup -- data encoding, our ring-CNOT ansatz
versus their tunable CRX/CRZ entanglers (Section~\ref{sec:method}), or
optimizer choice -- is the actual confound; our results cannot distinguish
between these. We are also careful with terminology here, since neither
prior paper's account of trainability loss for the circuit family closest
to ours uses "barren plateau": \citet{hqpinn2025flow} calls it "excessive
expressibility," based on final-loss curves rather than a direct gradient
measurement, and \citet{farea2025qcpinn} uses "barren plateau" explicitly
but only for its continuous-variable (qumode) circuits, a different
paradigm from the qubit-based ansatz both they and we build on. Our own
phase-2 gradient-norm measurement (Section~\ref{sec:experiments}) is, as
far as this literature review found, the only \emph{direct} circuit-only
gradient measurement across a qubit-count/depth sweep for a qubit-based
hybrid PINN. We offer this as a contribution, but state plainly, as
already noted in Section~\ref{sec:experiments}, that our own numbers do
not cleanly resolve the underlying question either.

One theoretical result changes how that non-resolution should be framed,
though: \citet{holmes2022expressibility} show that gradient magnitude and
ansatz expressibility are not independent, competing explanations but two
faces of the same relationship -- more expressive ansätze induce flatter
cost landscapes and correspondingly smaller typical gradients. Read
through that lens, \texttt{wide} having both the best accuracy \emph{and}
the smallest final gradient norm among your four quantum configurations is
not the contradiction it looks like under a naive "expressivity good,
vanishing gradient bad" framing -- it is closer to what
\citet{holmes2022expressibility}'s bound would predict for a more
expressive circuit that has not yet entered a true (exponential,
system-size-dependent) barren plateau: still trainable, but with a flatter
landscape than \texttt{narrow} or \texttt{baseline}. We stop short of
calling this a confirmed barren plateau, however: our qubit-count range
(2--6) is almost certainly too small for the asymptotic, exponential-decay
regime that the barren-plateau literature
\citep{mcclean2018barren,holmes2022expressibility} actually characterizes.
What our sweep shows is consistent with the small-system-size end of the
expressibility-gradient relationship those results describe -- a flatter
but still-trainable landscape as the circuit widens -- not evidence for or
against an eventual barren plateau at larger qubit counts, which our range
cannot distinguish from ordinary flattening.

\subsection{Limitations}
\begin{itemize}
  \item All quantum-circuit simulation is exact and noiseless
    (\texttt{default.qubit}); results say nothing about performance on
    real NISQ hardware, where shot noise and decoherence would likely
    widen rather than close the gap. (Farea et al.'s feasibility analysis
    of their Cross-mesh circuit under a stated noise model is the closest
    precedent for what a hardware-realistic version of this comparison
    would need to account for.)
  \item Both branches were run on CPU. This asymmetry is worth reading
    carefully rather than at face value: the classical branch could in
    practice be GPU-accelerated, which would shrink its training time
    further, while the quantum branch's cost is dominated by classically
    simulating the circuit and would not benefit comparably from the same
    hardware -- so a GPU-enabled comparison would likely widen, not narrow,
    the wall-clock gap reported here. Under future fault-tolerant quantum
    hardware this asymmetry could of course reverse, but that is a
    statement about hardware that does not yet exist, not about the
    noiseless-simulator comparison in this paper. We also note that
    neither \citet{farea2025qcpinn} nor \citet{hqpinn2025flow} give full
    hardware specifications sufficient for a precise cross-paper timing
    comparison -- a field-wide reporting gap, not a limitation specific to
    our study.
  \item Single PDE family (linear diffusion, 3 BC types); results may not
    transfer to nonlinear or higher-dimensional PDEs, or to the
    discontinuous/shock-type solutions where \citet{hqpinn2025flow}
    independently find classical dominance for different reasons.
  \item The two ablations use 3 seeds per (scenario, configuration), but
    the headline classical-vs-quantum comparison in
    Table~\ref{tab:main-results} uses a single seed per scenario -- it is
    the only result in the paper reported as a point estimate rather than
    a mean $\pm$ std, and the largest-magnitude claims (the 53--89$\times$
    wall-clock ratios) come from that single-seed table. This is a real
    limitation rather than a disclosure formality: the ablations show
    run-to-run standard deviations on the same order as some of the
    effects being compared, so we cannot yet rule out that the main
    comparison's specific ratios would shift, within their same broad
    range, under a different seed. Multi-seed replication of
    Table~\ref{tab:main-results} is the clearest concrete next step before
    this result is submission-ready.
  \item Phase-2 ablation varies qubit count and depth but not entangler
    type (fixed CNOT throughout) or data-encoding scheme --
    \citet{farea2025qcpinn} suggest entangler tunability (CRX/CRZ vs.\
    fixed CNOT) affects convergence reliability in their own qubit-circuit
    sweep, which we do not test and flag as the clearest concrete next
    ablation.
\end{itemize}


\section{Conclusion}\label{sec:conclusion}
\todo{One paragraph, no new citations. Suggested shape: (1) restate RQ1/RQ2
and the one-line answer -- quantum branch underperforms on both accuracy
and wall-clock cost across all six scenarios, and the gap survives
hyperparameter tuning; (2) state the diagnostic finding from ablation 2
honestly -- circuit structure matters but the effect does not cleanly
match either an expressivity-limited or barren-plateau account; (3) end on
the methodological point from the Discussion (wall-clock accounting) as
the actionable takeaway for the field, since that is the contribution most
likely to generalize beyond this specific PDE.}


\bibliographystyle{plainnat}
\bibliography{references}

\end{document}
